\documentclass[11pt]{article}

\usepackage[utf8]{inputenc}
\usepackage[T1]{fontenc}

\usepackage[
	a4paper,
	left = 2cm,
	right = 2cm,
	bottom = 1cm, includefoot,
	top = 1.2cm, includehead,
]{geometry}

\usepackage{amsmath,amssymb,amsfonts}
\RequirePackage{physics}
\usepackage{siunitx}
\usepackage[version=4]{mhchem}
\usepackage{chemfig}
\AtBeginDocument{\RenewCommandCopy\qty\SI}
\AtBeginDocument{\RenewCommandCopy\ce\chem}

\usepackage{enumitem}
\usepackage{booktabs}

\usepackage{hyperref}

\title{\textbf{01 -- Atomic Structure and Chemical Bonding}}
\author{Rodrigo Sánchez-Martínez}
\date{\today}

% --
\begin{document}
	
\maketitle	
	
%.
\tableofcontents

\section{Introduction}

\section{Electronic Structure of the Atom}
	\subsection{The nuclear atom}
	\subsection{Wave--particle duality and the Schrödinger equation}
	\subsection{Atomic orbitals: s, p, d, f}
	\subsection{Electron spin}

\section{Quantum Numbers and Atomic Orbitals}
	\subsection{\texorpdfstring{Principal quantum number $n$}{Principal quantum number n}}
	\subsection{\texorpdfstring{Angular momentum quantum number $l$}{Angular momentum quantum number l}}
	\subsection{\texorpdfstring{Magnetic quantum number $m_l$}{Magnetic quantum number m sub l}}
	\subsection{\texorpdfstring{Spin quantum number $m_s$}{Spin quantum number m sub s}}
	\subsection{Orbital shapes and radial/angular nodes}

\section{Electron Configuration and the Periodic Table}
	\subsection{Aufbau principle}
	\subsection{Pauli exclusion principle}
	\subsection{Hund's rule}
	\subsection{Electron configuration notation}
	\subsection{Valence electrons and core electrons}
	\subsection{Periodic table structure: blocks, groups, periods}

\section{Periodic Trends}
	\subsection{Atomic and ionic radius}
	\subsection{Ionization energy}
	\subsection{Electron affinity}
	\subsection{Electronegativity}

\section{Ionic Bonding}
	\subsection{Electron transfer and formation of ions}
	\subsection{Lattice energy}
	\subsection{Born--Haber cycle (brief)}
	\subsection{Properties of ionic compounds}

\section{Covalent Bonding}
	\subsection{Electron sharing and bond formation}
	\subsection{\texorpdfstring{Sigma ($\sigma$) and pi ($\pi$) bonds}{Sigma and pi bonds}}
	\subsection{Bond order and bond multiplicity}
	\subsection{Polar vs.\ nonpolar covalent bonds}

\section{Coordinate (Dative) Bonding}
	\subsection{Definition and formation}
	\subsection{Lewis acids and Lewis bases}
	\subsection{Examples in coordination chemistry}

\section{Metallic Bonding}
	\subsection{Electron-sea model}
	\subsection{Band theory: an introductory picture}
	\subsection{Conductors, semiconductors, and insulators}
	\subsection{Properties of metals}

\section{Bond Polarity and Electronegativity Difference}
	\subsection{Dipole moment of a bond}
	\subsection{Electronegativity scales (Pauling, Mulliken)}
	\subsection{Predicting bond character from electronegativity difference}

\section{Hybridization}
	\subsection{\texorpdfstring{$\mathrm{sp}$ hybridization}{sp hybridization}}
	\subsection{\texorpdfstring{$\mathrm{sp^2}$ hybridization}{sp2 hybridization}}
	\subsection{\texorpdfstring{$\mathrm{sp^3}$ hybridization}{sp3 hybridization}}
	\subsection{\texorpdfstring{$\mathrm{sp^3d}$ and $\mathrm{sp^3d^2}$ (brief)}{sp3d and sp3d2 (brief)}}

\section{Molecular Orbital Theory: A Qualitative Introduction}
	\subsection{Bonding and antibonding orbitals}
	\subsection{MO diagrams for homonuclear diatomic molecules}
	\subsection{MO theory vs.\ valence bond theory}

\section{Intermolecular Forces}
	\subsection{Van der Waals (London dispersion) forces}
	\subsection{Dipole--dipole interactions}
	\subsection{Hydrogen bonding}
	\subsection{Relevance to layered and van der Waals materials}

\section{Bond Energy, Bond Length, and Bond Order}
	\subsection{Definitions and typical magnitudes}
	\subsection{Trends across bond order}
	\subsection{Connection to vibrational spectroscopy}

\section{Summary}

\section{Key Terminology}

\section{Worked Examples}

\section{Exercises}

\end{document}
	
\end{document}